\section{Nichtfunktionale Anforderungen}
Neben den funktionalen und technischen Anforderungen spielen nichtfunktionale Eigenschaften eine wesentliche Rolle für die Qualität des Systems.
Die Zuverlässigkeit des Systems ist insbesondere im Dauerbetrieb von zentraler Bedeutung.
Gleichzeitig wird eine modulare Software Architektur angestrebt, um den Code wiederverwendbar und strukturiert aufzubauen, was die Wartbarkeit sicherstellen soll.
Alle nichtfunktionalen Anforderungen sind in Tabelle~\ref{tab:Nichtfunktionale Anforderung} aufgelistet.

\begin{table}[H]
\caption{Nichtfunktionale Anforderungen}
\centering
\begin{tabular}{p{0.5cm} p{10cm}}
\textbf{ID} & \textbf{Beschreibung} \\
\hline
N1 & Hohe Zuverlässigkeit im Dauerbetrieb \\
N2 & Energieeffizienter Betrieb trotz hoher LED-Anzahl \\
N3 & Wartbarkeit und klare Struktur des Codes \\
N4 & Ansprechendes und integrierbares Design \\
N5 & Benutzerfreundliche Bedienbarkeit \\
\end{tabular}
\label{tab:Nichtfunktionale Anforderung}
\end{table}